\chapter*{Resumen}
\addcontentsline{toc}{chapter}{Resumen}
\markboth{RESUMEN}{RESUMEN}

\begin{spacing}{1.5}

% BORRADOR: revisar al cerrar los resultados pendientes.
Este trabajo presenta un \textit{pipeline} \textit{on-premise} basado en modelos de lenguaje pequeños (\textit{Small Language Models}, SLM) para identificar errores en los logs de la Dirección de Tecnologías de la Información y Comunicación (DTIC) de la Universidad de Cuenca, notificar al desarrollador y recomendar acciones correctivas. Se utilizó el corpus del microservicio \textit{ms-users} (309 archivos, 27{,}752{,}836 eventos), que un prefiltro reduce a 223{,}104 eventos relevantes agrupados en 4{,}463 \textit{batches} de 50 eventos. Se construyó un \textit{ground truth} por \textit{batch} que combina la votación de tres modelos de lenguaje de gran tamaño (\textit{Large Language Models}, LLM) ($\kappa$ de Fleiss $=0.886$), la revisión de 329 \textit{batches} por dos expertos, reglas derivadas de su criterio y propagación por similitud, con 130 \textit{batches} de error (2.91\,\%). El consenso de los tres LLM sobreestimó el error: solo el 12.5\,\% de los \textit{batches} con voto de error unánime lo eran según los expertos. Se compararon ocho SLM con evaluación semántica mediante \textit{embeddings}. Contra la referencia por consenso de LLM, granite-4.1-3b obtuvo los mejores resultados ($F_1=0.951$ y un área bajo la curva [\textit{Area Under the Curve}, AUC] de 0.997); contra el \textit{ground truth} corregido por expertos, todos los modelos mostraron \textit{recall} alto (hasta 0.984) y baja precisión, y Qwen3.5-9B alcanzó el mejor $F_1$ (0.295). El \textit{pipeline} integra notificación a Slack, recomendador local, inyección de errores y \textit{dashboard}; los experimentos de inyección de errores quedan pendientes.

\end{spacing}

\hspace{1.27cm}\textit{Palabras clave del autor:} modelos de lenguaje pequeños, análisis de logs, detección de anomalías, observabilidad, \textit{on-premise}, \textit{ground truth}

\vfill % empuja el bloque CC hacia el pie de página

\begin{minipage}[t]{0.5\linewidth}
  \includegraphics[width=5cm]{figs/logos/image.png}
\end{minipage}
\par\vspace{0.25cm}
\begin{spacing}{0.5}
\begin{footnotesize}
El contenido de esta obra corresponde al derecho de expresión de los autores y no compromete el pensamiento institucional de la Universidad de Cuenca ni desata su responsabilidad frente a terceros. Los autores asumen la responsabilidad por la propiedad intelectual y los derechos de autor.
\par\vspace{0.6em} % una línea en blanco de separación
\textbf{Repositorio Institucional:} \href{https://dspace.ucuenca.edu.ec/}{\textcolor{blue}{\underline{\nolinkurl{https://dspace.ucuenca.edu.ec/}}}}
\end{footnotesize}
\end{spacing}
\vspace{0.45cm}\null % pequeño margen entre el enlace y el pie de página
